% Risk Studio -- slides for the CO4031 presentation (XeLaTeX).
% Build:  latexmk -xelatex slides.tex
% Layout mirrors the HTML deck: all coordinates are CSS px on a 1920x1080
% canvas, origin at the top-left corner, y downwards (see the theme file).
% Speaker notes: the note command at the end of each frame; README.md
% explains how to print them.
\documentclass[aspectratio=169,12pt]{beamer}
\usetheme{riskstudio}
% \setbeameroption{show notes on second screen=right}
% \setbeameroption{show only notes}

\title{Risk Studio – Thuyết trình CO4031}
\author{Doãn Phương Hùng Cường, Nguyễn Thiên Hào, Nguyễn Khắc Khánh Huy}

\begin{document}

% ===================================================================== 1
\begin{frame}{Risk Studio}
	\rsslide[dark]{}{%
		\rsT[text=rsEyebrowC]{128}{128}{\fEyebrow}%
		{\MakeUppercase{CO4031 · Kho dữ liệu và Hệ hỗ trợ quyết định · HK261}}
		\rsGradeBars{128}{362.7}
		\rsT{128}{402.7}{\fBig}{Risk Studio}
		\rsTw[text=rsTealSoft]{128}{542.7}{1200}{\fHthreeReg}%
		{Kho dữ liệu và hệ hỗ trợ quyết định thẩm định tín dụng}
		\rsT[text=rsTealSoft]{128}{788.2}{\fSmallS}{Lớp L02 · Nhóm sinh viên}
		\rsT{128}{826.6}{\fBody}{Doãn Phương Hùng Cường · 2310381}
		\rsT{128}{870.4}{\fBody}{Nguyễn Thiên Hào · 2310853}
		\rsT{128}{914.2}{\fBody}{Nguyễn Khắc Khánh Huy · 2311186}
		\rsT[anchor=north east]{1792}{875.8}{\fBodyS}{GVHD: ThS. Đỗ Thanh Thái}
		\rsT[anchor=north east,text=rsMutedC]{1792}{919.6}{\fSmall}%
		{Dữ liệu: Home Credit Default Risk (Kaggle)}
	}
	\note{%
		Hùng Cường – khoảng 20 giây. Chào thầy và các bạn. Nhóm chúng em gồm ba sinh viên lớp
		L02: Doãn Phương Hùng Cường, Nguyễn Thiên Hào và Nguyễn Khắc Khánh Huy, dưới sự hướng
		dẫn của thầy Đỗ Thanh Thái. Đề tài là Risk Studio: kho dữ liệu và hệ hỗ trợ quyết định
		thẩm định tín dụng trên bộ dữ liệu Home Credit của Kaggle. Mỗi bạn trình bày khoảng năm
		phút: em trình bày bài toán và kho dữ liệu, Thiên Hào trình bày phân tích và thang hạng,
		Khánh Huy trình bày kết quả chính sách, demo và hạn chế.}
\end{frame}

% ===================================================================== 2
\begin{frame}{Bài toán thẩm định tín dụng}
	\rsslide{Hùng Cường}{%
		\rsHead{01 · Bối cảnh}{Bài toán thẩm định tín dụng}
		% left column (860 px)
		\rsT[text=rsTeal]{128}{278.4}{\fBig}{356.255}
		\rsT{128}{402.4}{\fBodyS}{hồ sơ vay cần một đề xuất}
		\rsTw[text=rsMuted]{128}{452.2}{860}{\fSmall}%
		{307.511 khoản vay lịch sử (đã có kết quả trả nợ) + 48.744 hồ sơ mới (cần thẩm định)}
		% right column: three decisions
		\rsT[text=rsMuted]{1052}{278.4}{\fSmallS}{Mỗi hồ sơ nhận một trong ba đề xuất}
		\rsPill[256]{1052}{326.8}{approve}{Approve}
		\rsT{1328}{330.1}{\fBody}{Duyệt}
		\rsPill[256]{1052}{387.2}{review}{Manual review}
		\rsT{1328}{390.5}{\fBody}{Thẩm định thêm}
		\rsPill[256]{1052}{447.6}{reject}{Reject}
		\rsT{1328}{450.9}{\fBody}{Từ chối}
		% three questions, pinned to the bottom
		\rsCard{128}{734}{533.33}{186}
		\rsCard{693.33}{734}{533.33}{186}
		\rsCard{1258.67}{734}{533.33}{186}
		\rsT[text=rsTeal]{161}{767}{\fSmallS}{Q2 · Phân khúc}
		\rsTw{161}{811.4}{467.33}{\fBody}{Phân khúc nào rủi ro cao và chiếm bao nhiêu danh mục?}
		\rsT[text=rsTeal]{726.33}{767}{\fSmallS}{Q4 · Từng hồ sơ}
		\rsTw{726.33}{811.4}{467.33}{\fBody}{Hồ sơ này nên xử lý ra sao, và vì sao?}
		\rsT[text=rsTeal]{1291.67}{767}{\fSmallS}{Q5 · Chính sách}
		\rsTw{1291.67}{811.4}{467.33}{\fBody}{Đổi ngưỡng thì tỷ lệ duyệt và tỷ lệ vỡ nợ đổi ra sao?}
	}
	\note{%
		Hùng Cường – khoảng 50 giây. Home Credit cho vay tiêu dùng với nhóm khách hàng có ít
		hoặc không có lịch sử tín dụng. Mỗi hồ sơ là một quyết định cân bằng hai rủi ro: duyệt
		quá lỏng thì mất vốn, duyệt quá chặt thì mất khách tốt. Thẩm định thủ công lại thiếu
		nhất quán và khó giải thích. Hệ thống của nhóm đưa ra một trong ba đề xuất: duyệt, thẩm
		định thêm hoặc từ chối, cho 356.255 hồ sơ, gồm 307.511 khoản vay lịch sử đã biết kết quả
		và 48.744 hồ sơ mới. Ba câu hỏi chính nhóm muốn trả lời là: phân khúc nào rủi ro cao,
		một hồ sơ cụ thể nên xử lý ra sao và vì sao, và nếu đổi chính sách thì tỷ lệ duyệt và tỷ
		lệ vỡ nợ thay đổi thế nào.}
\end{frame}

% ===================================================================== 3
\begin{frame}{Người dùng và cấp quyết định}
	\rsslide{Hùng Cường}{%
		\rsHead{02 · Người dùng}{Người dùng và cấp quyết định}
		\rsTable{128}{278.4}{\fBody}{20}{465.92,449.28,748.8}{l,l,l}
		\rsRow[th]{61.8}{Người dùng & Cấp quyết định & Quyết định cần hỗ trợ}
		\rsRow{61.8}{Trưởng phòng rủi ro & Chiến lược, chiến thuật & Đặt ngưỡng hạng, luật và chính sách}
		\rsRow{61.8}{Chuyên viên thẩm định & Tác nghiệp & Duyệt, thẩm định thêm hay từ chối hồ sơ}
		\rsRow{61.8}{Ban điều hành & Giám sát & Theo dõi KPI của danh mục}
		% Simon model, pinned to the bottom
		\rsT[text=rsMuted]{128}{646.6}{\fSmallS}%
		{Mô hình Simon: mỗi giai đoạn ra quyết định có trang web riêng}
		\rsCard{128}{695}{512}{225}
		\rsCard{704}{695}{512}{225}
		\rsCard{1280}{695}{512}{225}
		\draw[rs conn] (640,807.5) -- (704,807.5);
		\draw[rs conn] (1216,807.5) -- (1280,807.5);
		\rsT[text=rsTeal]{161}{728}{\fSmallS}{Nhận thức · Intelligence}
		\rsT{161}{772.4}{\fBodyS}{Overview · Analytics · Explorer}
		\rsTw[text=rsMuted]{161}{822.2}{446}{\fSmall}{Phát hiện phân khúc và yếu tố rủi ro}
		\rsT[text=rsTeal]{737}{728}{\fSmallS}{Thiết kế · Design}
		\rsT{737}{772.4}{\fBodyS}{Playground}
		\rsTw[text=rsMuted]{737}{822.2}{446}{\fSmall}{Thử ngưỡng, luật và chính sách theo hạng}
		\rsT[text=rsTeal]{1313}{728}{\fSmallS}{Lựa chọn · Choice}
		\rsT{1313}{772.4}{\fBodyS}{Application details}
		\rsTw[text=rsMuted]{1313}{822.2}{446}{\fSmall}{Chọn đề xuất cho từng hồ sơ, kèm lý do}
	}
	\note{%
		Hùng Cường – khoảng 50 giây. Hệ thống phục vụ ba nhóm người dùng. Trưởng phòng rủi ro ra
		quyết định chiến lược và chiến thuật: đặt ngưỡng hạng, các luật và chính sách cho từng
		hạng. Chuyên viên thẩm định ra quyết định tác nghiệp trên từng hồ sơ. Ban điều hành giám
		sát qua các KPI. Thẩm định tín dụng là quyết định bán cấu trúc: một phần mô hình hoá
		được bằng dữ liệu, phần còn lại cần con người. Vì vậy hệ thống chỉ đề xuất, không tự
		quyết. Theo mô hình Simon, giai đoạn nhận thức dùng Overview, Analytics và Explorer;
		giai đoạn thiết kế dùng Playground; giai đoạn lựa chọn dùng trang Application details.
		Tiếp theo, em xin trình bày phần kho dữ liệu.}
\end{frame}

% ===================================================================== 4
\begin{frame}{Dữ liệu nguồn}
	\rsslide{Hùng Cường}{%
		\rsHead{03 · Kho dữ liệu}{Dữ liệu nguồn}
		\rsT[text=rsTeal]{128}{278.4}{\fBig}{8}
		\rsT[text=rsMuted]{128}{394.4}{\fSmall}{bảng CSV từ Kaggle}
		\rsT[text=rsTeal]{128}{446.8}{\fBig}{2,7 GB}
		\rsT[text=rsMuted]{128}{562.8}{\fSmall}{dung lượng CSV}
		\rsT[text=rsTeal]{128}{615.2}{\fBig}{58,5}
		\rsT[text=rsMuted]{128}{731.2}{\fSmall}{triệu dòng (khoảng)}
		\rsTable{632}{278.4}{\fSmall}{20}{464,696}{l,l}
		\rsRow[th]{56.4}{Bảng & Mỗi dòng là}
		\rsRow{56.4}{application\_train & Hồ sơ vay có nhãn TARGET (lịch sử)}
		\rsRow{56.4}{application\_test & Hồ sơ vay chưa có nhãn (hồ sơ mới)}
		\rsRow{56.4}{bureau & Khoản vay tại tổ chức khác (Credit Bureau)}
		\rsRow{56.4}{bureau\_balance & Quá hạn theo tháng của khoản bureau}
		\rsRow{56.4}{previous\_application & Hồ sơ vay cũ tại Home Credit}
		\rsRow{56.4}{installments\_payments & Một kỳ trả góp: phải trả so với thực trả}
		\rsRow{56.4}{POS\_CASH\_balance & Số dư, quá hạn theo tháng (POS, tiền mặt)}
		\rsRow{56.4}{credit\_card\_balance & Số dư, hạn mức, quá hạn của thẻ tín dụng}
		% simulated-dates call-out with a clock icon
		\rsCardSoft{128}{826}{1664}{130.2}
		\draw[rsTealDeep,line width=4\rspx,line cap=round,line join=round]
		(184,891.1) circle (20) (184,879.1) -- (184,891.1) -- (192,895.1);
		\rsT{232}{854}{\fBodyS}{Thêm một nguồn MÔ PHỎNG: ngày nộp hồ sơ và ngày có kết quả}
		\rsT[text=rsMuted]{232}{895.8}{\fSmall}%
		{Dữ liệu gốc chỉ có số ngày tính lùi, nên ngày được sinh bằng hcdw synth-dates.}
	}
	\note{%
		Hùng Cường – khoảng 40 giây. Dữ liệu nguồn là bộ Home Credit Default Risk trên Kaggle: 8
		bảng CSV, khoảng 2,7 GB và khoảng 58,5 triệu dòng. Bảng application là trung tâm; tập
		train có nhãn TARGET nên đóng vai khoản vay lịch sử, tập test chưa có nhãn nên đóng vai
		hồ sơ mới. Các bảng còn lại là lịch sử tín dụng: bureau, hồ sơ cũ, trả góp, POS và thẻ.
		Một điểm quan trọng: dữ liệu gốc không có ngày thật, chỉ có số ngày tính lùi từ lúc nộp
		hồ sơ. Vì vậy nhóm thêm một nguồn mô phỏng cho ngày nộp và ngày có kết quả. Mọi biểu đồ
		theo thời gian trong phần sau đều dựa trên ngày mô phỏng này, nên không đọc thành xu
		hướng nghiệp vụ.}
\end{frame}

% ===================================================================== 5
\begin{frame}{Kiến trúc tổng thể}
	\rsslide{Hùng Cường}{%
		\rsHead{04 · Kho dữ liệu}{Kiến trúc tổng thể}
		\begin{scope}[shift={(128,278.4)}]
			\fill[rsBand,rounded corners=16\rspx] (0,0) rectangle ++(1664,352);
			\fill[rsBand,rounded corners=16\rspx] (0,384) rectangle ++(1664,256);
			\rsT[text=rsEyebrowC]{32}{16}{\fEyebrow}{\MakeUppercase{Offline · CLI hcdw}}
			\rsT[text=rsEyebrowC]{32}{400}{\fEyebrow}{\MakeUppercase{Serving · hcdw serve}}
			\rsNode{plain}{32}{64}{304}{136}{Kaggle CSV}{8 bảng · 2,7 GB}
			\rsNode{plain}{464}{64}{304}{136}{ETL · hcdw etl}{làm sạch, gộp, nạp, 50 kiểm tra}
			\rsNode{key}{896}{64}{304}{136}{DuckDB}{staging · dw · mart}
			\rsNode{plain}{1328}{64}{304}{136}{Train + score}{LightGBM \rsRA{} hạng A–E}
			\rsNode{sim}{32}{236}{400}{100}{Ngày mô phỏng}{application\_dates.csv}
			\rsNode{plain}{32}{448}{680}{136}{React SPA · Risk Studio}%
			{Overview · Analytics · Explorer · Application details · Playground}
			\rsNode{plain}{896}{448}{304}{136}{FastAPI}{/api/* · OLAP + DSS · chỉ đọc kho}
			\draw[rs conn] (336,132) -- (464,132);
			\draw[rs conn] (768,132) -- (896,132);
			\draw[rs conn both] (1200,132) -- (1328,132);
			\draw[rs conn,rs dashed] (432,286) -| (616,200);
			\draw[rs conn] (1048,200) -- (1048,448);
			\draw[rs conn] (1480,200) |- (1200,516);
			\draw[rs conn] (712,516) -- (896,516);
			\rsT[text=rsMuted,anchor=south]{1264}{124}{\fSmall}{đọc / ghi}
			\rsT[text=rsMuted]{1064}{280}{\fSmall}{chỉ đọc}
			\rsT[text=rsMuted]{1496}{280}{\fSmall}{đọc mô hình}
			\rsT[text=rsMuted]{728}{468}{\fSmall}{HTTP JSON}
		\end{scope}
	}
	\note{%
		Hùng Cường – khoảng 50 giây. Hệ thống chia hai phần. Phần offline chạy bằng dòng lệnh
		hcdw: ETL đọc 8 file CSV của Kaggle và nguồn ngày mô phỏng, làm sạch, gộp, nạp vào một
		file DuckDB gồm ba lược đồ staging, dw và mart. Sau đó bước train học mô hình LightGBM
		ngay trên kho, và bước score chấm mọi hồ sơ rồi ghi kết quả hạng A đến E trở lại kho.
		Phần serving gồm API FastAPI chỉ đọc kho và mô hình đã lưu, và giao diện React có năm
		trang chính. Trình duyệt không bao giờ gửi SQL; API kiểm tra tên chiều và chỉ số theo
		danh sách cho phép. Offline và serving dùng chung mã chấm điểm, nên chấm hàng loạt và
		chấm một hồ sơ cho cùng kết quả.}
\end{frame}

% ===================================================================== 6
\begin{frame}{Mô hình đa chiều}
	\rsslide{Hùng Cường}{%
		\rsHead{05 · Kho dữ liệu}{Mô hình đa chiều}
		\rsT[text=rsTeal]{128}{278.4}{\fBig}{4}
		\rsT{128}{394.4}{\fBody}{bảng sự kiện (fact)}
		\rsT[text=rsTeal]{128}{464.2}{\fBig}{8}
		\rsT{128}{580.2}{\fBody}{bảng chiều (dimension)}
		\rsTw[text=rsMuted]{128}{650}{340}{\fSmall}%
		{Star schema, không snowflake: mỗi chiều chỉ cần một phép join}
		% analytical part: fact constellation
		\rsCard{516}{278.4}{755.69}{504.6}
		\rsT{549}{311.4}{\fBodyS}{Kho phân tích · fact constellation}
		\fill[rsTeal,rounded corners=10\rspx] (549,365.2) rectangle ++(336.85,86.8);
		\fill[rsTeal,rounded corners=10\rspx] (901.85,365.2) rectangle ++(336.85,86.8);
		\rsT[text=rsPage]{565}{375.2}{\fSmallS}{fact\_loan\_application}
		\rsT[text=rsTealSoft]{565}{409.6}{\fSmall}{1 dòng = 1 hồ sơ vay}
		\rsT[text=rsPage]{917.85}{375.2}{\fSmallS}{fact\_bureau\_credit}
		\rsT[text=rsTealSoft]{917.85}{409.6}{\fSmall}{1 dòng = 1 khoản bureau}
		\foreach \x/\y/\t in {549/468/dim\_customer, 901.85/468/dim\_time,
				549/540.4/dim\_credit\_history, 901.85/540.4/dim\_contract,
				549/612.8/dim\_region, 901.85/612.8/dim\_bureau\_category}{
				\rsBox{\x}{\y}{336.85}{56.4}{10}{none}{rsNodeLine}{2}
				\rsT{\x+18}{\y+12}{\fSmall}{\t}}
		\rsTw[text=rsMuted]{549}{685.2}{689.69}{\fSmall}%
		{dim\_time: role-playing · dim\_credit\_history: mini-dimension · dim\_bureau\_category: junk}
		% DSS part
		\rsCard{1319.69}{278.4}{472.31}{454.2}
		\rsT{1352.69}{311.4}{\fBodyS}{Phần DSS}
		\fill[rsTeal,rounded corners=10\rspx] (1352.69,365.2) rectangle ++(406.31,86.8);
		\rsT[text=rsPage]{1368.69}{375.2}{\fSmallS}{fact\_risk\_assessment}
		\rsT[text=rsTealSoft]{1368.69}{409.6}{\fSmall}{1:1 với hồ sơ vay}
		\fill[rsTeal,rounded corners=10\rspx] (1352.69,468) rectangle ++(406.31,86.8);
		\rsT[text=rsPage]{1368.69}{478}{\fSmallS}{fact\_rule\_hit}
		\rsT[text=rsTealSoft]{1368.69}{512.4}{\fSmall}{1 dòng = 1 luật vi phạm}
		\rsBox{1352.69}{570.8}{406.31}{56.4}{10}{none}{rsNodeLine}{2}
		\rsT{1370.69}{582.8}{\fSmall}{dim\_risk\_grade · A–E}
		\rsBox{1352.69}{643.2}{406.31}{56.4}{10}{none}{rsNodeLine}{2}
		\rsT{1370.69}{655.2}{\fSmall}{dim\_rule · 6 luật}
		\rsTw[text=rsMuted]{128}{844.4}{1664}{\fBody}%
		{Hai fact chính dùng chung dim\_customer và dim\_time, nên phân tích chéo (drill-across) được từ hồ sơ vay sang khoản vay bureau.}
	}
	\note{%
		Hùng Cường – khoảng 45 giây. Kho dữ liệu gồm 4 bảng sự kiện và 8 bảng chiều. Phần phân
		tích là một fact constellation: fact\_loan\_application có hạt là một hồ sơ vay,
		fact\_bureau\_credit có hạt là một khoản vay ở bureau. Hai fact dùng chung chiều khách
		hàng và chiều thời gian, nên có thể phân tích chéo. dim\_time đóng nhiều vai, ví dụ ngày
		nộp và ngày có kết quả; dim\_credit\_history là mini-dimension gộp hàng chục triệu dòng
		lịch sử; dim\_bureau\_category là junk dimension. Phần DSS thêm fact\_risk\_assessment,
		một dòng cho mỗi hồ sơ đã chấm, và fact\_rule\_hit, một dòng cho mỗi luật bị vi phạm,
		cùng hai chiều hạng và luật nạp từ file cấu hình.}
\end{frame}

% ===================================================================== 7
\begin{frame}{ETL và cổng chất lượng}
	\rsslide{Hùng Cường}{%
		\rsHead{06 · Kho dữ liệu}{ETL và cổng chất lượng}
		\begin{scope}[shift={(128,278.4)}]
			\rsNode{plain}{0}{0}{320}{150}{1 · Extract}{8 CSV + ngày mô phỏng vào staging}
			\rsNode{plain}{448}{0}{320}{150}{2 · Clean}{Quy tắc R1–R10 (pandas)}
			\rsNode{plain}{896}{0}{320}{150}{3 · Aggregate}{Lịch sử \rsRA{} 1 dòng mỗi khách (SQL)}
			\rsNode{plain}{1344}{0}{320}{150}{4 · Load}{Dim, fact và view mart}
			\rsNode{gate}{1344}{240}{320}{150}{5 · Quality gate}{50 kiểm tra: PK, FK, số dòng…}
			\rsNode{ok}{896}{240}{320}{150}{Publish}{Thay hc.duckdb bằng file tạm}
			\rsNode{bad}{1344}{480}{320}{150}{Rollback}{Xoá file tạm, giữ kho cũ, exit 1}
			\rsT[text=rsTeal]{0}{240}{\fBig}{\textasciitilde 15 giây}
			\rsT[text=rsMuted]{0}{360}{\fSmall}{hcdw etl trên toàn bộ 2,7 GB CSV}
			\rsT{0}{424.4}{\fNsixty}{\textasciitilde 1,5 phút}
			\rsT[text=rsMuted]{0}{498.4}{\fSmall}{hcdw etl -\kern0pt-score: thêm train và chấm điểm}
			\draw[rs conn] (320,75) -- (448,75);
			\draw[rs conn] (768,75) -- (896,75);
			\draw[rs conn] (1216,75) -- (1344,75);
			\draw[rs conn] (1504,150) -- (1504,240);
			\draw[rs conn,draw=rsGA] (1344,315) -- (1216,315);
			\draw[rs conn,draw=rsGE] (1504,390) -- (1504,480);
			\rsT[text=rsGA]{1240}{263}{\fSmallS}{Đạt}
			\rsT[text=rsGE]{1520}{419}{\fSmallS}{Không đạt}
		\end{scope}
	}
	\note{%
		Hùng Cường – khoảng 45 giây. Quy trình ETL đi qua các bước extract, clean, aggregate và
		load. Extract đọc 8 file CSV và nguồn ngày mô phỏng vào staging; clean áp dụng mười quy
		tắc làm sạch R1 đến R10; aggregate gộp lịch sử tín dụng về một dòng mỗi khách; load nạp
		các bảng chiều, bảng sự kiện và view phân tích. Mọi bước ghi vào một file tạm. Sau đó
		cổng chất lượng chạy 50 kiểm tra như khoá chính, khoá ngoại, số dòng và đối soát tổng
		tiền. Nếu tất cả đạt, file tạm thay kho đang dùng; nếu có kiểm tra không đạt, file tạm
		bị xoá, kho cũ giữ nguyên. Toàn bộ ETL mất khoảng 15 giây, và khoảng 1,5 phút nếu chạy
		thêm train và chấm điểm. Tiếp theo, Thiên Hào sẽ trình bày phần phân tích.}
\end{frame}

% ===================================================================== 8
\begin{frame}{KPI chính}
	\rsslide{Thiên Hào}{%
		\rsHead{07 · Phân tích}{KPI chính}
		% 3 x 2 grid, columns 533.33 px, gap 32
		\rsCard{128}{278.4}{533.33}{284.2}
		\rsPill{161}{311.4}{scopeall}{All applications}
		\rsT{161}{367.8}{\fBig}{51,86\%}
		\rsT{161}{491.8}{\fBody}{Tỷ lệ duyệt · KPI-04}
		\rsCard{693.33}{278.4}{533.33}{284.2}
		\rsPill{726.33}{311.4}{scopeall}{All applications}
		\rsT{726.33}{367.8}{\fBig}{39,54\%}
		\rsT{726.33}{491.8}{\fBody}{Tỷ lệ chuyển thẩm định · KPI-05}
		\rsCard{1258.67}{278.4}{533.33}{284.2}
		\rsPill{1291.67}{311.4}{scopehist}{Historical loans}
		\rsT{1291.67}{367.8}{\fBig}{8,07\%}
		\rsT{1291.67}{491.8}{\fBody}{Tỷ lệ vỡ nợ · KPI-01}
		\rsCard{128}{594.6}{533.33}{284.2}
		\rsPill{161}{627.6}{scopeall}{All applications}
		\rsT[name=kpiA]{161}{684}{\fBig}{209,4}
		\node[anchor=base west] at ([xshift=10\rspx]kpiA.base east) {\fHthree tỷ};
		\rsT{161}{808}{\fBody}{Dư nợ cấp mới · KPI-02}
		\rsCard{693.33}{594.6}{533.33}{284.2}
		\rsPill{726.33}{627.6}{scopeall}{All applications}
		\rsT[name=kpiB]{726.33}{684}{\fBig}{15,76}
		\node[anchor=base west] at ([xshift=10\rspx]kpiB.base east) {\fHthree tỷ};
		\rsT{726.33}{808}{\fBody}{Dư nợ chịu rủi ro · KPI-03}
		\rsCardSoft{1258.67}{594.6}{533.33}{284.2}
		\rsT[text=rsTeal]{1290.67}{625.5}{\fNsixty}{356.255}
		\rsTw{1290.67}{699.5}{469.33}{\fBody}%
		{hồ sơ đã chấm = 307.511 khoản vay lịch sử + 48.744 hồ sơ mới}
		\rsTw[text=rsMuted]{1290.67}{783.1}{469.33}{\fSmall}%
		{Tỷ lệ vỡ nợ chỉ tính được trên khoản vay lịch sử. Tiền giữ đơn vị gốc.}
	}
	\note{%
		Thiên Hào – khoảng 1 phút. Đây là năm KPI chính trên trang Overview. Mỗi ô ghi rõ phạm
		vi. Trên mọi hồ sơ đã chấm, chính sách hiện hành duyệt 51,86\% và chuyển thẩm định thêm
		39,54\%. Dư nợ cấp mới là 209,4 tỷ, và dư nợ chịu rủi ro, tức tổng tiền vay nhân xác
		suất vỡ nợ, là 15,76 tỷ, giữ nguyên đơn vị tiền trong dữ liệu gốc. Riêng tỷ lệ vỡ nợ
		8,07\% chỉ tính trên khoản vay lịch sử, vì chỉ chúng có kết quả trả nợ. Hai phạm vi này
		có nhãn All applications và Historical loans trên giao diện, để người xem không so nhầm
		hai con số khác mẫu số.}
\end{frame}

% ===================================================================== 9
\begin{frame}{OLAP trên trang Analytics}
	\rsslide{Thiên Hào}{%
		\rsHead{08 · Phân tích}{OLAP trên trang Analytics}
		\rsTable{128}{278.4}{\fBody}{20}{276,644}{l,l}
		\rsRow[th]{61.8}{Phép OLAP & Thao tác trên Analytics}
		\rsRow{61.8}{Slice & Một chip lọc: cố định một chiều}
		\rsRow{61.8}{Dice & Nhiều chip lọc trên từ hai chiều}
		\rsRow{61.8}{Roll-up & Interval lên Year; \rsUp{} hoặc × bỏ chiều}
		\rsRow{61.8}{Drill-down & Interval xuống Month; \rsDown{} hoặc + chiều}
		\rsRow{61.8}{Pivot & \rsSwap{} đổi chiều hàng và chiều cột}
		\rsRow{61.8}{Drill-through & Bấm một ô \rsRA{} mở Explorer}
		\rsCardSoft{128}{788.4}{920}{131.6}
		\rsTw{160}{816.4}{856}{\fBody}%
		{Tùy thao tác, trang gửi 1, 3 hoặc 5 yêu cầu tới API; mỗi yêu cầu chạy một câu SQL GROUPING SETS (15–60 ms).}
		\rsShot{1100}{128}{692}{792}{images/page_analytics.png}
	}
	\note{%
		Thiên Hào – khoảng 1 phút. Trang Analytics minh hoạ các phép OLAP trên cùng một khối dữ
		liệu. Slice là thêm một chip lọc cố định một chiều; dice là nhiều chip trên từ hai
		chiều. Roll-up và drill-down là đổi Interval giữa năm, quý, tháng, hoặc dùng nút mũi
		tên, thêm và bỏ chiều trên bảng Cube. Pivot là nút đổi hàng và cột. Bấm một ô sẽ mở
		Explorer với đúng nhóm đó, gọi là drill-through. Tùy thao tác trên bảng Cube, Interval
		hay bộ lọc, trang gửi 1, 3 hoặc 5 yêu cầu tới API; mỗi yêu cầu chạy một câu SQL
		GROUPING SETS trong khoảng 15 đến 60 mili giây. Lưu ý: trục thời gian ở đây dùng ngày
		mô phỏng.}
\end{frame}

% ===================================================================== 10
\begin{frame}{Insight từ dữ liệu}
	\rsslide{Thiên Hào}{%
		\rsHead{09 · Phân tích}{Insight từ dữ liệu}
		% left: default rate by credit history
		\rsCard{128}{278.4}{808}{619.6}
		\rsT{161}{311.4}{\fHthree}{Theo lịch sử tín dụng}
		\rsT[text=rsMuted]{161}{370.6}{\fSmall}%
		{Tỷ lệ vỡ nợ · khoản vay lịch sử · đường đứt: chung 8,07\%}
		\foreach \x/\h/\c/\v in {207/116/rsTeal/{10,12\%}, 443/88/rsTeal/{7,62\%},
				679/185/rsTealDeep/{16,11\%}}{
				\fill[\c] (\x,791) [rounded corners=8\rspx] -- (\x,791-\h) -- (\x+180,791-\h)
				[sharp corners] -- (\x+180,791) -- cycle;
				\node[anchor=north] at (\x+90,791-\h-45.8) {\fBodyS\rsstrut\v\rsstrut};}
		\fill[rsNodeLine] (161,791) rectangle ++(744,2);
		\draw[rsMuted,line width=2\rspx,dash pattern=on 6\rspx off 6\rspx]
		(161,699) -- ++(744,0);
		\node[anchor=north] at (297,809) {\fSmall\rsstrut No bureau\rsstrut};
		\node[anchor=north] at (533,809) {\fSmall\rsstrut Clean history\rsstrut};
		\node[anchor=north] at (769,809) {\fSmall\rsstrut Current overdue\rsstrut};
		% right: default rate by final grade, same scale
		\rsCard{984}{278.4}{808}{619.6}
		\rsT{1017}{311.4}{\fHthree}{Theo hạng cuối: tăng đều A \rsRA{} E}
		\rsT[text=rsMuted]{1017}{370.6}{\fSmall}%
		{Tỷ lệ vỡ nợ · khoản vay lịch sử · cùng thang với biểu đồ trái}
		\foreach \x/\h/\r/\c/\v/\g in {1057/7/4/rsGA/{0,62\%}/A, 1195/41/8/rsGB/{3,54\%}/B,
				1333/95/8/rsGC/{8,29\%}/C, 1471/140/8/rsGD/{12,15\%}/D, 1609/322/8/rsGE/{28,02\%}/E}{
				\fill[\c] (\x,791) [rounded corners=\r\rspx] -- (\x,791-\h) -- (\x+112,791-\h)
				[sharp corners] -- (\x+112,791) -- cycle;
				\node[anchor=north] at (\x+56,791-\h-45.8) {\fBodyS\rsstrut\v\rsstrut};
				\rsChip{\x+28}{809}{\c}{\g}}
		\fill[rsNodeLine] (1017,791) rectangle ++(744,2);
	}
	\note{%
		Thiên Hào – khoảng 1 phút. Hai insight chính. Thứ nhất, lịch sử tín dụng ảnh hưởng rõ
		đến khả năng trả nợ: khách đang quá hạn ở bureau vỡ nợ 16,11\%, gấp đôi mức chung
		8,07\%; khách không có dữ liệu bureau vỡ nợ 10,12\%, còn khách có lịch sử sạch chỉ
		7,62\%. Đây là lý do lịch sử tín dụng là thuộc tính đầu tiên của luật phân khúc. Thứ
		hai, thang hạng cuối phân tách rủi ro tốt: tỷ lệ vỡ nợ thực tế tăng đều từ 0,62\% ở hạng
		A lên 28,02\% ở hạng E, đúng yêu cầu BR-07. Hai biểu đồ dùng chung một thang đo. Tiếp
		theo, em xin trình bày cách hệ thống xếp hạng.}
\end{frame}

% ===================================================================== 11
\begin{frame}{Thang hạng A–E}
	\rsslide{Thiên Hào}{%
		\rsHead{10 · Hỗ trợ quyết định}{Thang hạng A–E}
		\begin{scope}[shift={(128,278.4)}]
			\rsNode{plain}{0}{82}{240}{136}{Hồ sơ}{dữ liệu trong kho}
			\rsNode{plain}{400}{0}{380}{136}{1 · Hạng điểm}{PD của mô hình \rsRA{} ngưỡng A–E}
			\rsNode{plain}{400}{164}{380}{136}{2 · Hạng luật}{Hạng xấu nhất trong các luật vi phạm}
			\rsNode{key}{940}{82}{380}{136}{3 · Hạng cuối}{Hạng XẤU HƠN của hai hạng}
			\rsNode{plain}{1420}{82}{244}{136}{4 · Chính sách}{Quyết định + hạn mức}
			\draw[rs conn] (240,150) -- (320,150) -- (320,68) -- (400,68);
			\draw[rs conn] (240,150) -- (320,150) -- (320,232) -- (400,232);
			\draw[rs conn] (780,68) -- (860,68) -- (860,150) -- (940,150);
			\draw[rs conn] (780,232) -- (860,232) -- (860,150) -- (940,150);
			\draw[rs conn] (1320,150) -- (1420,150);
		\end{scope}
		\rsT[text=rsMuted]{128}{664.8}{\fSmallS}%
		{Một thang chung cho hạng điểm và hạng luật; mỗi hạng gắn đúng một chính sách}
		\foreach \i/\c/\g/\t/\p/\d/\m in {
				0/rsGA/A/{< 1\%}/approve/Approve/{Tối đa 8× thu nhập},
				1/rsGB/B/{1\% – 7\%}/approve/Approve/{Tối đa 6× thu nhập},
				2/rsGC/C/{7\% – 13\%}/review/Manual review/{Tối đa 4× thu nhập},
				3/rsGD/D/{13\% – 20\%}/review/Manual review/{Tối đa 3× thu nhập},
				4/rsGE/E/{≥ 20\%}/reject/Reject/{Không cấp hạn mức}}{
				\pgfmathsetmacro\cx{128+337.6*\i}
				\rsCard{\cx}{713.2}{313.6}{206.8}
				\rsChip{\cx+29}{738.2}{\c}{\g}
				\rsT{\cx+97}{747.3}{\fBodyS}{\t}
				\rsPill{\cx+29}{806.2}{\p}{\d}
				\rsT[text=rsMuted]{\cx+29}{862.6}{\fSmall}{\m}}
	}
	\note{%
		Thiên Hào – khoảng 1 phút 10 giây. Hệ hỗ trợ quyết định dùng một thang hạng A đến E duy
		nhất, xác định theo xác suất vỡ nợ với các ngưỡng 1, 7, 13 và 20\%. Mỗi hồ sơ được đánh
		giá từ hai nguồn. Hạng điểm đến từ PD do mô hình ước lượng. Hạng luật là hạng xấu nhất
		trong các luật mà hồ sơ vi phạm; không vi phạm luật nào thì không có hạng luật. Hạng
		cuối là hạng xấu hơn trong hai hạng, nên luật chỉ hạ hạng, không bao giờ nâng. Mỗi hạng
		gắn một chính sách: A và B duyệt, hạn mức tối đa 8 và 6 lần thu nhập; C và D chuyển thẩm
		định thêm với 4 và 3 lần; E từ chối.}
\end{frame}

% ===================================================================== 12
\begin{frame}{Mô hình PD và luật}
	\rsslide{Thiên Hào}{%
		\rsHead{11 · Hỗ trợ quyết định}{Mô hình PD và luật}
		\rsT[text=rsTeal]{128}{278.4}{\fBig}{0,7586}
		\rsT{128}{398.4}{\fBody}{AUC out-of-fold}
		\rsT[text=rsTeal]{128}{468.2}{\fBig}{0,0678}
		\rsT{128}{588.2}{\fBody}{Brier, PD đã hiệu chỉnh}
		\rsCardSoft{128}{658}{520}{161.2}
		\rsTw{160}{686}{456}{\fSmallS}{LightGBM 5-fold OOF + hiệu chỉnh isotonic}
		\rsTw[text=rsMuted]{160}{758.8}{456}{\fSmall}{26 đặc trưng, không dùng giới tính}
		\rsTable{712}{278.4}{\fBody}{20}{237.6,518.4,129.6,194.4}{l,l,l,r}
		\rsRow[th]{61.8}{Mã & Luật & Hạng & Vỡ nợ}
		\rsRow{61.8}{BR-SEG-01 & Phân khúc rủi ro cao & \textcolor{rsGD}{D, E} & 15,12\%}
		\rsRow{61.8}{BR-RF-01 & Hồ sơ cũ bị từ chối nhiều & \textcolor{rsGD}{D} & 15,10\%}
		\rsRow{61.8}{BR-RF-02 & Trả trễ hơn 20\% kỳ & \textcolor{rsGC}{C} & 11,99\%}
		\rsRow{61.8}{BR-RF-03 & Từng quá hạn hơn 90 ngày & \textcolor{rsGC}{C} & 9,83\%}
		\rsRow{61.8}{BR-RF-04 & Đang quá hạn ở bureau & \textcolor{rsGD}{D} & 16,11\%}
		\rsRow{61.8}{BR-RF-05 & Trả trễ hơn 50\% kỳ & \textcolor{rsGD}{D} & 16,44\%}
		\rsTw[text=rsMuted]{712}{727}{1080}{\fSmall}%
		{Vỡ nợ: tỷ lệ vỡ nợ lịch sử của nhóm vi phạm, quy về thang A–E thành hạng của luật.}
	}
	\note{%
		Thiên Hào – khoảng 50 giây. Mô hình PD là LightGBM với kiểm định chéo 5 phần. Mỗi khoản
		vay lịch sử nhận PD từ một mô hình chưa từng thấy nó, sau đó PD được hiệu chỉnh bằng
		isotonic để khớp tỷ lệ vỡ nợ thật, vì ngưỡng hạng là xác suất. AUC đạt 0,7586 và Brier
		0,0678. Mô hình không dùng giới tính. Bên phải là sáu luật: một luật phân khúc BR-SEG-01
		và năm luật nghiệp vụ BR-RF-01 đến 05. Hạng của mỗi luật lấy từ tỷ lệ vỡ nợ lịch sử của
		chính nhóm vi phạm, ví dụ khách đang quá hạn ở bureau vỡ nợ 16,11\% nên luật có hạng D.
		Mỗi luật có mã và điều kiện rõ ràng, nên chuyên viên kiểm tra lại được. Tiếp theo, Khánh
		Huy sẽ trình bày kết quả của chính sách.}
\end{frame}

% ===================================================================== 13
\begin{frame}{Kết quả chính sách}
	\rsslide[accent]{Khánh Huy}{%
		\rsHead{12 · Hỗ trợ quyết định}{Kết quả chính sách}
		\fill[rsPage,rounded corners=16\rspx] (128,278.4) rectangle ++(1664,256.8);
		\rsT[text=rsInk]{160}{310.4}{\fSmallS}{Đề xuất trên mọi hồ sơ đã chấm (356.255)}
		\begin{scope}
			\clip[rounded corners=12\rspx] (160,358.8) rectangle ++(1600,96);
			\fill[rsApprove] (160,358.8) rectangle ++(829.84,96);
			\fill[rsReview] (989.84,358.8) rectangle ++(632.70,96);
			\fill[rsReject] (1622.54,358.8) rectangle ++(137.46,96);
		\end{scope}
		\rsT[anchor=north,text=rsPaper]{574.92}{387.9}{\fBodyS}{Approve 51,86\%}
		\rsT[anchor=north,text=rsInk]{1306.19}{387.9}{\fBodyS}{Manual review 39,54\%}
		\rsT[anchor=north,text=rsPaper]{1691.27}{387.9}{\fBodyS}{8,59\%}
		\rsT[text=rsMuted]{160}{470.8}{\fSmall}%
		{Approve = hạng A, B · Manual review = hạng C, D · Reject 8,59\% = hạng E}
		\rsT{128}{720.4}{\fBig}{3,37\%}
		\rsTw[text=rsTealSoft]{128}{844.4}{760}{\fBody}%
		{vỡ nợ của nhóm được duyệt, so với 8,07\% chung (khoản vay lịch sử)}
		\rsT{984}{720.4}{\fBig}{13,55\%}
		\rsTw[text=rsTealSoft]{984}{844.4}{640}{\fBody}%
		{hồ sơ bị luật hạ hạng so với hạng điểm (mọi hồ sơ)}
	}
	\note{%
		Khánh Huy – khoảng 50 giây. Với chính sách mặc định, trên mọi hồ sơ đã chấm, hệ thống đề
		xuất duyệt 51,86\%, chuyển thẩm định thêm 39,54\% và từ chối 8,59\%. Kết quả quan trọng
		nhất: nhóm được duyệt chỉ vỡ nợ 3,37\%, so với 8,07\% của toàn bộ khoản vay lịch sử, tức
		chưa bằng một nửa. Luật đã hạ hạng 13,55\% hồ sơ so với hạng điểm. Khối lượng thẩm định
		thêm gần 40\% là lựa chọn thận trọng; đây cũng là điểm nhóm sẽ quay lại ở phần hạn chế.
		Câu hỏi tiếp theo là: nếu đổi chính sách thì sao? Đó là việc của trang Playground.}
\end{frame}

% ===================================================================== 14
\begin{frame}{Mô phỏng chính sách trên Playground}
	\rsslide{Khánh Huy}{%
		\rsHead{13 · Hỗ trợ quyết định}{Mô phỏng chính sách trên Playground}
		\rsTable{128}{278.4}{\fSmall}{20}{179.2,286.72,197.12,232.96}{l,l,r,r}
		\rsRow[th]{88.8}{Preset & Thay đổi & Approval rate & Default rate (approved)}
		\rsRow[hl]{56.4}{Hiện hành & PD 1/7/13/20\% & 51,86\% & 3,37\%}
		\rsRow{56.4}{Approve C & Hạng C \rsRA{} Approve & 74,90\% & 4,87\%}
		\rsRow{56.4}{Loosen & PD 2/9/15/25\% & 59,10\% & 3,85\%}
		\rsRow{56.4}{Tighten & PD 0,5/4/10/15\% & 29,31\% & 2,22\%}
		\rsRow{56.4}{Rules off & Chỉ dùng hạng điểm & 61,74\% & 3,52\%}
		\rsCardSoft{128}{673.2}{896}{134.2}
		\rsTw{160}{701.2}{832}{\fBodyS}{Duyệt nhiều hơn hiện hành thì vỡ nợ nhóm duyệt tăng}
		\rsTw[text=rsMuted]{160}{747}{832}{\fSmall}%
		{Mô phỏng dùng lại PD và lần vi phạm đã lưu, khoảng 37 ms mỗi lần.}
		\rsShot{1072}{278.4}{720}{640}{images/page_playground.png}
	}
	\note{%
		Khánh Huy – khoảng 50 giây. Trang Playground cho trưởng phòng rủi ro thử chính sách
		trước khi áp dụng. Có thể đổi ngưỡng hạng, quyết định và hạn mức của từng hạng, bật tắt
		từng luật. Bảng bên trái là bốn preset so với hiện hành. Chuyển hạng C sang duyệt nâng
		tỷ lệ duyệt lên 74,90\% nhưng vỡ nợ nhóm duyệt tăng lên 4,87\%. Siết ngưỡng làm vỡ nợ
		giảm còn 2,22\% nhưng chỉ duyệt 29,31\%. Tắt luật thì duyệt 61,74\% và vỡ nợ 3,52\%.
		Không có phương án nào vừa duyệt nhiều vừa ít vỡ nợ, nên Playground giúp người quyết
		định thấy rõ đánh đổi. Tiếp theo, em xin demo trực tiếp trên web.}
\end{frame}

% ===================================================================== 15
\begin{frame}{Demo trực tiếp}
	\rsslide{Khánh Huy}{%
		\rsHead{14 · Demo}{Demo trực tiếp}
		% five demo steps; card 4 (the worked case) has a 2 px teal border
		\foreach \y/\h/\n/\t/\s in {
				278.4/106.2/1/Overview/{5 KPI, ô waffle và Segment analysis},
				396.6/106.2/2/Analytics/{Slice, roll-up và pivot trên bảng Cube},
				514.8/106.2/3/Explorer/{Lọc một nhóm hồ sơ, so với mọi hồ sơ},
				753.2/106.2/5/Playground/{Thử preset Approve C, Reset rồi Tighten}}{
				\rsCard{128}{\y}{760}{\h}
				\fill[rsTeal] (181,\y+53.1) circle (24);
				\node[anchor=base,text=rsPaper] at (181,\y+63.39) {\fBodyS\n};
				\rsT{229}{\y+17}{\fBodyS}{\t}
				\rsT[text=rsMuted]{229}{\y+56.8}{\fSmall}{\s}}
		\rsBox{128}{633}{760}{108.2}{16}{rsPaper}{rsTeal}{2}
		\fill[rsTeal] (182,687.1) circle (24);
		\node[anchor=base,text=rsPaper] at (182,697.39) {\fBodyS 4};
		\rsT{230}{651}{\fBodyS}{Application details · 257416}
		\rsT[text=rsMuted]{230}{690.8}{\fSmall}{Hạng điểm C, bị luật hạ xuống D: Manual review}
		\rsShot{936}{278.4}{856}{600}{images/page_application_details.png}
		\rsT[text=rsMuted]{936}{890.4}{\fSmall}{Ngày trên giao diện là ngày mô phỏng.}
	}
	\note{%
		Khánh Huy – khoảng 2 phút 20 giây, chuyển sang trình duyệt. Một: Overview, chỉ năm KPI
		và biểu đồ tỷ lệ vỡ nợ theo lịch sử tín dụng. Hai: Analytics, thêm một chip lọc, đổi
		Interval sang Year để roll-up, rồi bấm nút đổi hàng cột. Ba: Explorer, lọc một nhóm và
		so với mọi hồ sơ. Bốn: mở Application details, hồ sơ 257416. Hạng điểm là C, nhưng hồ sơ
		vi phạm ba luật, trong đó có luật hạng D, nên hạng cuối là D và đề xuất là thẩm định
		thêm; bên trái là các yếu tố SHAP. Năm: Playground, bấm Approve C, bấm Reset rồi bấm
		Tighten để thấy các ô KPI đổi (vì các nút preset áp dụng chồng lên cấu hình đang mở).
		Lưu ý ngày trên giao diện là ngày mô phỏng. Nếu mạng hoặc máy có sự cố, dùng ảnh chụp
		trên slide này.}
\end{frame}

% ===================================================================== 16
\begin{frame}{Hạn chế và hướng phát triển}
	\rsslide{Khánh Huy}{%
		\rsHead{15 · Đánh giá}{Hạn chế và hướng phát triển}
		\rsT[text=rsMuted]{128}{278.4}{\fSmallS}{Hạn chế}
		\rsT[text=rsTeal]{992}{278.4}{\fSmallS}{Hướng phát triển}
		\foreach \y/\a/\b/\c/\d in {
				324.8/{Ngày nộp và ngày kết quả là mô phỏng}/{Xu hướng theo thời gian không có ý nghĩa nghiệp vụ}%
				/{Dùng ngày thật, nạp tăng dần}/{Incremental load theo ngày, thay vì dựng lại toàn bộ},
				447/{Hiệu chỉnh PD chưa kiểm định độc lập}/{Isotonic được khớp và đo trên cùng PD OOF}%
				/{Tách tập holdout}/{Khớp trên phần còn lại, đo Brier trên holdout},
				569.2/{Luật đổi nhiều, cải thiện ít}/{Duyệt 61,74\% \rsRA{} 51,86\%; vỡ nợ duyệt 3,52\% \rsRA{} 3,37\%}%
				/{Đánh giá lại từng luật}/{Tắt hoặc đổi ngưỡng, đo tác động trên Playground},
				691.4/{Thẩm định thêm 39,54\% hồ sơ}/{Nhiều việc cho chuyên viên thẩm định}%
				/{Tìm nhóm hạng C rủi ro thấp}/{Dùng Explorer và Playground, thử chính sách riêng},
				813.6/{What-if chưa đơn điệu}/{Mô hình cây: vài hồ sơ đi ngược hướng nghiệp vụ}%
				/{Thêm monotone constraints}/{Cho thu nhập và số tiền vay trong LightGBM}}{
				\rsCard{128}{\y}{800}{108.2}
				\rsT{157}{\y+17}{\fBodyS}{\a}
				\rsT[text=rsMuted]{157}{\y+58.8}{\fSmall}{\b}
				\draw[rs conn] (928,\y+54.1) -- (992,\y+54.1);
				\rsCardSoft{992}{\y}{800}{108.2}
				\rsT{1020}{\y+16}{\fBodyS}{\c}
				\rsT[text=rsMuted]{1020}{\y+57.8}{\fSmall}{\d}}
	}
	\note{%
		Khánh Huy – khoảng 45 giây. Nhóm nhận thấy năm hạn chế chính, mỗi hạn chế có một hướng
		khắc phục. Một, ngày là dữ liệu mô phỏng, nên cần ngày thật và nạp tăng dần. Hai, bộ
		hiệu chỉnh PD được khớp và đo trên cùng dữ liệu, nên cần một tập holdout để kiểm định
		thật. Ba, luật làm tỷ lệ duyệt giảm gần 10 điểm phần trăm nhưng vỡ nợ nhóm duyệt chỉ
		giảm khoảng 0,15 điểm, nên cần đánh giá lại từng luật trên Playground. Bốn, gần 40\% hồ
		sơ phải thẩm định thêm; có thể tìm nhóm hạng C rủi ro thấp để tự động hoá. Năm, phân
		tích what-if chưa đơn điệu; có thể thêm ràng buộc đơn điệu cho mô hình.}
\end{frame}

% ===================================================================== 17
\begin{frame}{Cảm ơn · Hỏi đáp}
	\rsslide[dark]{}{%
		\rsGradeBars{128}{392.2}
		\rsT{128}{436.2}{\fBig}{Cảm ơn · Hỏi đáp}
		\rsT[text=rsTealSoft]{128}{580.2}{\fHthreeReg}%
		{Doãn Phương Hùng Cường · Nguyễn Thiên Hào · Nguyễn Khắc Khánh Huy}
		\rsT[text=rsMutedC]{128}{655.4}{\fSmall}%
		{Risk Studio · CO4031 HK261 · Lớp L02 · GVHD: ThS. Đỗ Thanh Thái}
	}
	\note{%
		Khánh Huy – khoảng 15 giây. Phần trình bày của nhóm đến đây là kết thúc. Tóm lại, nhóm
		đã xây dựng một kho dữ liệu từ 8 bảng Kaggle, một hệ hỗ trợ quyết định xếp hạng A đến E
		kèm giải thích, và một trang mô phỏng chính sách. Cảm ơn thầy và các bạn đã lắng nghe;
		nhóm rất mong nhận được câu hỏi và góp ý. Gợi ý phân công khi trả lời, theo bảng phân
		công trong báo cáo: yêu cầu nghiệp vụ, KPI và hệ hỗ trợ quyết định (thang hạng, mô hình,
		luật, chính sách, Playground) do Thiên Hào; kho dữ liệu, ETL, OLAP, backend và giao diện
		web do Khánh Huy; đánh giá kết quả, hạn chế, hướng phát triển và báo cáo do Hùng Cường.}
\end{frame}

\end{document}
